\section{Related Work}

%% TODO: 
%% * add multi-agent system literature. Use this survey as an entry point https://ieeexplore.ieee.org/stamp/stamp.jsp?tp=&arnumber=8352646
%% * discuss motivation of developing multi-agent collaboration in GenAI era, cite earlier work such as ChatDev, MetaGPT, MAD, 

\paragraph{Multi-agent GenAI Systems} An agent is defined as ``an entity which is placed in an environment and
senses different parameters that are used to make a decision
based on the goal of the entity.''~\cite{kapoor2024ai}. The motivation for multi-agent systems is to have agents collaborate on a complex task that could not have been accomplished by a single agent. With the rise of LLMs, researchers have proposed leveraging the reasoning and planning capabilities of these models to develop more sophisticated multi-agent systems. Examples of such multi-agent LLM systems include MetaGPT~\cite{hong2024metagpt} and CAMEL~\cite{li2023camel}.  MetaGPT is one of the first multi-agent LLM projects that tries to mimic a software company. Developers can provide a standard operating procedure and MetaGPT tries to assign roles to various agents.
CAMEL, or Communicative Agent Framework, promotes independent collaboration between LLM agents. Its key innovation is the use of "inception prompting," a method that steers conversational agents to complete tasks. Beyond its practical applications, CAMEL also functions as a research platform. It enables the creation and analysis of conversational data, offering valuable insights into the behavior and interactions of communicative agents. Across these works, much of the emphasis is on customization and coordination, which was much less prominent in traditional multi-agent systems~\cite{singh1994multiagent,sycara1998multiagent}.

\paragraph{Multi-agent Frameworks and Platforms} CrewAI~\cite{crewai} is designed to enhance task execution by organizing agents into specialized roles, similar to team members in a crew. This approach emphasizes task decomposition, where a complex problem is broken down into smaller, manageable subtasks. Each agent is assigned a specific role based on its expertise, allowing for efficient problem-solving. AutoGen~\cite{wu2024autogen} represents a significant advancement in enabling multi-agent systems to engage in sophisticated conversations. This framework allows agents to communicate and collaborate by sharing information and refining their outputs through iterative interactions. By simulating human-like dialogues, AutoGen enables agents to negotiate, plan, and execute tasks collaboratively. 

LangGraph~\cite{langgraph} introduces an innovative framework for organizing agent interactions using directed acyclic graphs (DAGs). This structure allows for clear visualization and management of dependencies between tasks and agents. By leveraging DAGs, LangGraph optimizes the flow of information and decision-making processes among agents. This approach enhances the system’s ability to handle complex, interdependent tasks by ensuring that each agent’s actions are informed by the outcomes of preceding steps. 

%Vertex AI integrates multi-agent collaboration concepts into a comprehensive platform that supports the development, deployment, and scaling of AI models. By incorporating tools for orchestrating agent interactions and managing workflows, Vertex AI provides a robust infrastructure for building collaborative AI systems.

\paragraph{End-to-end Agent Evaluation Frameworks}

While our work on assertion-based benchmarking is ongoing, there are other works in the literature that are similar. AgentEval~\cite{arabzadeh-etal-2024-assessing} has a multi-agent setup where there are three agents to evaluate conversation trajectories: 1) CriticAgent, 2) QuantifierAgent, 3) VerifierAgent. The CriticAgent takes the task description as input and outputs a list of criteria for task success. The QuantifierAgent then assesses whether the trajectories meet the criteria and the outputs here can be a scalar value. The VerifierAgent will finally verify that the evaluation is accurate and complete. Note that the criteria proposed by the CriticAgent is much more coarse-grained like “clarity” and “efficiency”.

AgentQuest~\cite{gioacchini-etal-2024-agentquest} proposes to measure success with a “progress rate” that is based on a set of milestones. The progress rate measures how many milestones are completed and milestones are defined to be “environment hidden states the agent needs to reach to get the final solution of the task”. AgentQuest focuses benchmarking agents for game-like datasets like ALFWorld and Sudoku, so state naturally comes with the environment when the agent plays the game. These milestones can either be externally annotated or programmatically defined within the simulation.

Most recently, ToolSandbox~\cite{lu2024toolsandbox} is the work in literature most similar to our evaluation setup. ToolSandbox also includes a user simulator and environment executor for tools. Before the simulations, they also pre-define “milestones” and “minefields” for each session. Milestones are events that must occur during the conversation and minefields are those that should not happen. These milestones and minefields are similar to our assertions that cover actions and agent behaviors. In addition, their implementation stores an “Execution Context” that contains the “world state” to mimic tasks that manipulate a resource like a database.
